\chapter{Derivation of the Jacobian of the Re-Parametrization}
\label{app:calc_dydx}

Recall the re-parameterization
\begin{equation*}
\mathbf{y} = \bigl[ \lambda_1,\dots,\lambda_{n-1},\,s,\,\gamma_1,\dots,\gamma_n\bigr]
\in \mathbb{R}^{2n},
\qquad
\mathbf{x} = g(\mathbf{y}) = \bigl[ L_1 , \dots , L_n, \, c_1 , \dots , c_n \bigr],
\end{equation*}
with
\begin{equation*}
c_i = e^{\gamma_i},
\qquad
\rho = \frac{1}{1+e^{-s}},
\qquad
L_i = L\,\rho\,\alpha_i,
\end{equation*}
and $\alpha_i$ defined by
\begin{equation*}
\alpha_i
=\begin{cases}
\dfrac{e^{\lambda_i}}{1+\sum_{k=1}^{n-1} e^{\lambda_k}}, & i\in\{1,2,\dots,n-1\},\\[16pt]
\dfrac{1}{1+\sum_{k=1}^{n-1} e^{\lambda_k}}, & i=n.
\end{cases}
\end{equation*}

For $i\in\{1,2,\dots,n\}$ and $j\in\{1,2,\dots,n-1\}$:

\subsection*{Derivatives of $\alpha_i$}
Let
\begin{equation}
D = 1 + \sum_{k=1}^{n-1} e^{\lambda_k},
\end{equation}
and note that
\begin{equation}
\frac{\partial D}{\partial \lambda_j} = e^{\lambda_j}.
\end{equation}

For $i\neq n$, $\alpha_i = \dfrac{e^{\lambda_i}}{D}$. Using the quotient rule,

\begin{align*}
\frac{\partial \alpha_i}{\partial \lambda_j}
&=
\dfrac{D \frac{\partial}{\partial \lambda_j}(e^{\lambda_i}) - e^{\lambda_i} \frac{\partial}{\partial \lambda_j}(D)}{D^2}
=
\frac{
\delta_{ij} e^{\lambda_i} D -
e^{\lambda_i} e^{\lambda_j}
}{
D^2
}\\
&= \delta_{ij} \frac{ e^{\lambda_i}}{D} - \frac{ e^{\lambda_i}}{D} \frac{ e^{\lambda_j}}{D} = \delta_{ij} \alpha_i - \alpha_i \alpha_j,
\end{align*}
where $\delta_{ij}$ is the kronecker delta function.

For $i = n$, $\alpha_n = \dfrac{1}{D}$, so
\begin{equation}
\frac{\partial \alpha_n}{\partial \lambda_j}
=
-\frac{1}{D^2} \frac{\partial D}{\partial \lambda_j}
=
-\frac{e^{\lambda_j}}{D^2}
=
-\frac{1}{D}\frac{e^{\lambda_j}}{D}
=
-\alpha_n \alpha_j.
\end{equation}

To summarize,
\begin{equation}
\frac{\partial \alpha_i}{\partial \lambda_j}
=
\begin{cases}
\alpha_i(1-\alpha_i), & i=j,\\[6pt]
-\alpha_i \alpha_j, & i\neq j.
\end{cases}
\end{equation}

\subsection*{Derivative of $\rho$}

\begin{equation*}
\frac{\partial \rho}{\partial s}
=
\rho(1-\rho)
\end{equation*}

\subsection*{Derivatives of $L_i$}

\begin{align*}
\frac{\partial L_i}{\partial \lambda_j}
&=
\frac{\partial L_i}{\partial \alpha_i} \frac{\partial \alpha_i}{\partial \lambda_j} =
L\rho
\begin{cases}
\alpha_i(1-\alpha_i), & i=j\\[4pt]
-\alpha_i \alpha_j, & i\neq j
\end{cases}
\\[6pt]
\frac{\partial L_i}{\partial s}
&=
\frac{\partial L_i}{\partial \rho} \frac{\partial \rho}{\partial s} =
L\,\alpha_i\,\rho(1-\rho)
\\[6pt]
\frac{\partial L_i}{\partial \gamma_j}
&=0
\end{align*}

\subsection*{Derivatives of $c_i$}

\begin{align*}
\frac{\partial c_i}{\partial \gamma_j}
&=
\begin{cases}
c_i, & i=j\\[4pt]
0, & i\neq j
\end{cases}
\\[6pt]
\frac{\partial c_i}{\partial \lambda_j} &=0\\
\frac{\partial c_i}{\partial s} & = 0
\end{align*}

\subsection*{Final Structure of the Jacobian}

Collecting all terms, the Jacobian of the re-parametrization has the block structure
\begin{equation*}
\frac{\partial \mathbf{x}}{\partial \mathbf{y}}
=
\begin{bmatrix}
\displaystyle \frac{\partial \mathbf{L}}{\partial \boldsymbol{\lambda}} &
\displaystyle \frac{\partial \mathbf{L}}{\partial s} &
\mathbf{0}_{n\times n}
\\[8pt]
\mathbf{0}_{n\times (n-1)} &
\mathbf{0}_{n\times 1} &
\operatorname{diag}(c_1,\dots,c_n)
\end{bmatrix} \in \mathbb{R}^{2n\times 2n}.
\end{equation*}